\documentclass[12pt, a4paper]{article}

\usepackage[margin=1.8cm,top=2cm,bottom=2cm]{geometry}
\usepackage{amsmath,amssymb,amsfonts}
\usepackage{graphicx}
\usepackage[colorlinks=true,linkcolor=blue,citecolor=blue,urlcolor=blue]{hyperref}
\usepackage{bm,xcolor}
\usepackage[numbers,sort&compress]{natbib}
\usepackage{float}

\newcommand{\Geff}{G_{\mathrm{eff}}}
\newcommand{\cs}{c_s}
\newcommand{\meff}{m_{\mathrm{eff}}}
\newcommand{\heff}{\hbar_{\mathrm{eff}}}
\newcommand{\lP}{\ell_{\mathrm{P}}}
\newcommand{\mP}{m_{\mathrm{P}}}
\newcommand{\tP}{t_{\mathrm{P}}}

\begin{document}

\title{Emergent Time as Phase Dynamics\\
in the Logarithmic Superfluid Vacuum}

\author{B. B. Kulangiev\\\small{Haifa, Israel}}
\date{April 2026}
\maketitle

\begin{abstract}
\noindent
We show that time is not a fundamental dimension
in the logarithmic superfluid vacuum framework but
an emergent bookkeeping parameter generated by the
phase dynamics of the condensate order parameter.
The Madelung decomposition $\psi = \sqrt{\rho}\,
e^{iS/\heff}$ defines time operationally as the
phase accumulation rate: $dS/dt = \mu$, where the
chemical potential $\mu(x) = \meff\,c_s^2(x)/\heff$
is the local clock frequency. The condensate is a
three-dimensional object; its state at any instant
is fully specified by $\rho(\mathbf{x})$ and
$S(\mathbf{x})$. The ``time dimension'' perceived
by internal observers (phonons) is an artifact of
writing the phonon equation of motion as a wave
equation on a Lorentzian manifold. We demonstrate
that all temporal phenomena in general relativity
have condensate counterparts: gravitational time
dilation is the spatial variation of $\mu(x)$, the
twin paradox is path-dependent phase accumulation,
the arrow of time is the irreversibility of phase
advance in an open system, and the Hubble flow is
the spatial phase winding rate. At the healing
length $\xi \sim \lP$, the phase counter saturates,
providing a minimum resolvable time interval
$\Delta t_{\min} \sim \tP$. At acoustic horizons
($v = c_s$), the phase becomes locked and the
emergent clock stops---reproducing infinite
gravitational redshift without a metric singularity.
The framework predicts that time did not exist
before the condensate: the nucleation of the vacuum
from the Bulk reservoir is simultaneously the origin
of space, time, and the physical constants.
\end{abstract}


% ====================================================================
\section{Introduction}
\label{sec:intro}
% ====================================================================

The status of time in fundamental physics remains
deeply problematic. In Newtonian mechanics, time is
an absolute external parameter. In special relativity,
it becomes a coordinate entangled with space. In
general relativity, it is a dynamical variable---the
metric determines the rate at which clocks tick.
In canonical quantum gravity, the Wheeler-DeWitt
equation~\cite{dewitt_1967} is timeless: the
Hamiltonian constraint $\hat{H}\Psi = 0$ contains no
time derivative, and the recovery of time evolution
from the ``frozen formalism'' remains an open
problem~\cite{isham_1993,kuchar_1992}.

The question ``what is time?'' has motivated diverse
proposals: Barbour's timeless mechanics, where time
is change~\cite{barbour_1999}; the thermal time
hypothesis of Connes and Rovelli, where time emerges
from thermodynamic equilibrium~\cite{connes_1994};
and Volovik's observation that in superfluid helium,
time is emergent from the condensate
dynamics~\cite{volovik_2003}. In all these
approaches, time is not fundamental but derived.

In companion
papers~\cite{kulangiev_2026a,kulangiev_2026b,
kulangiev_2026c,kulangiev_2026d,kulangiev_2026e},
we developed a superfluid vacuum model in which
gravity, dark energy, dark matter, and the Planck
scale all emerge from a logarithmic condensate.
Throughout these papers, the time coordinate $t$
appeared as a parameter in the wave equation and
acoustic metric. In this paper, we examine the
nature of this parameter and argue that it is not
a fundamental dimension of reality but an emergent
quantity generated by the phase dynamics of the
order parameter.


% ====================================================================
\section{The Phase Counter}
\label{sec:phase}
% ====================================================================

\subsection{Time from the Madelung decomposition}

The vacuum condensate is described by a complex
order parameter:
\begin{equation}
    \psi(\mathbf{x}, t) = \sqrt{\rho(\mathbf{x}, t)}\;
    \exp\!\left(\frac{iS(\mathbf{x}, t)}{\heff}\right)\,.
\label{eq:madelung}
\end{equation}
The state of the condensate at any ``instant'' is
fully specified by two three-dimensional fields:
the density $\rho(\mathbf{x})$ and the phase action
$S(\mathbf{x})$. The parameter $t$ enters only
through the Hamilton-Jacobi equation~\cite{kulangiev_2026a}:
\begin{equation}
    \frac{\partial S}{\partial t} = -\mu(\mathbf{x})\,,
\label{eq:ds_dt}
\end{equation}
where $\mu(\mathbf{x})$ is the local chemical
potential. In the rest frame of the condensate:
\begin{equation}
    \mu(\mathbf{x}) =
    \frac{\meff\,c_s^2(\mathbf{x})}{\heff}\,,
\label{eq:mu}
\end{equation}
with $c_s$ the local sound speed. The order
parameter oscillates as $\psi \propto e^{-i\mu t}$,
and $t$ is the parameter counting these oscillations.

This is operationally identical to how a clock works:
a clock is a physical system that oscillates at a
known frequency, and ``one second'' is a fixed number
of cycles. Here, the oscillator is the vacuum
condensate itself, and the ``frequency'' is $\mu/\heff$.
Time is not a dimension the condensate lives in;
it is a label the condensate generates by oscillating.

\subsection{The 3+1 illusion}

The acoustic metric experienced by phonons
(photons) is~\cite{kulangiev_2026a,visser_1998}:
\begin{equation}
    ds^2 = \frac{\rho}{c_s}\left[
    -(c_s^2 - v^2)\,dt^2
    - 2\,\mathbf{v}\cdot d\mathbf{x}\,dt
    + d\mathbf{x}^2\right]\,.
\label{eq:acoustic_metric}
\end{equation}
This has four coordinates $(t, \mathbf{x})$ and a
Lorentzian signature $(-,+,+,+)$. A phonon
propagating on this metric experiences a 3+1
dimensional spacetime.

But the metric~(\ref{eq:acoustic_metric}) is
constructed entirely from three-dimensional
quantities: $\rho(\mathbf{x})$, $c_s(\mathbf{x})$,
and $\mathbf{v}(\mathbf{x})$. The ``time
dimension'' is an artifact of writing the phonon
equation of motion as a wave equation on a curved
manifold. The underlying reality is a 3D condensate
whose configuration evolves; $t$ is the index
labelling the sequence of configurations.

This is Volovik's insight applied to the vacuum:
in superfluid helium, the phonon sees a 3+1
dimensional acoustic spacetime, but the helium
itself is a 3D fluid in a laboratory with an
external clock. For the vacuum condensate, there is
no external clock---the phase oscillation is the
only available timer, and the ``time dimension''
is entirely self-generated~\cite{volovik_2003}.


% ====================================================================
\section{Gravitational Time Dilation as Chemical
Potential Variation}
\label{sec:dilation}
% ====================================================================

\subsection{The mechanism}

The local clock rate is set by the chemical
potential $\mu(\mathbf{x})$. Near a gravitating
mass, the Painlev\'{e}-Gullstrand flow modifies
the local vacuum density and sound speed. From the
acoustic metric~(\ref{eq:acoustic_metric}), the
proper time interval is:
\begin{equation}
    d\tau = \sqrt{-g_{00}}\,dt
    = \sqrt{\frac{\rho\,c_s}{1}}\,
    \sqrt{1 - v^2/c_s^2}\,dt\,.
\label{eq:proper_time}
\end{equation}
In the weak-field limit ($v \ll c_s$):
\begin{equation}
    \frac{d\tau}{dt} \approx 1 - \frac{v^2}{2c_s^2}
    = 1 - \frac{GM}{c^2 r}\,,
\label{eq:redshift}
\end{equation}
recovering the standard gravitational redshift.
The physical content is: the condensate phase
advances more slowly where the flow velocity is
larger (i.e., deeper in the gravitational well),
because the kinetic energy of the flow reduces the
energy available for phase oscillation.

\subsection{Numerical example: the solar surface}

At the solar surface, $v^2/c^2 \sim GM_\odot/(c^2
R_\odot) \sim 2 \times 10^{-6}$. The chemical
potential is reduced by one part in $5 \times 10^5$:
\begin{equation}
    \frac{\delta\mu}{\mu}
    = -\frac{GM_\odot}{c^2 R_\odot}
    \approx -2.12 \times 10^{-6}\,.
\end{equation}
A clock on the solar surface accumulates
$2.12 \times 10^{-6}$ fewer phase cycles per second
than a clock at infinity. This is the Pound-Rebka
effect~\cite{pound_1960}, here derived from the
condensate phase dynamics rather than from the
metric tensor.


% ====================================================================
\section{The Twin Paradox as Path-Dependent Phase
Accumulation}
\label{sec:twin}
% ====================================================================

The twin paradox has a transparent interpretation in
the phase counter picture. Two observers $A$ and $B$
start at the same spacetime point with synchronized
phases $S_A = S_B$. Observer $B$ travels along a
worldline $\gamma_B$ through regions of varying
$\mu(\mathbf{x})$ and flow velocity
$\mathbf{v}(\mathbf{x})$, while $A$ remains
stationary.

The total phase accumulated by each observer is:
\begin{equation}
    \Delta S_i = -\int_{\gamma_i} \mu\,d\tau_i\,,
    \qquad i = A, B\,.
\label{eq:phase_path}
\end{equation}
When the observers reunite, the phase difference is:
\begin{equation}
    \Delta S_A - \Delta S_B
    = -\int_{\gamma_A}\!\mu\,d\tau_A
    + \int_{\gamma_B}\!\mu\,d\tau_B\,.
\label{eq:twin_phase}
\end{equation}
The observer who accumulated fewer phase cycles has
aged less. This is not a coordinate artifact: the
phase difference is a scalar quantity, measurable
by interfering the two condensate paths. Time
dilation is the physical fact that different
trajectories through the condensate accumulate
different numbers of phase oscillations.


% ====================================================================
\section{The Arrow of Time}
\label{sec:arrow}
% ====================================================================

\subsection{Phase advance is irreversible}

The condensate phase evolves as
$S(t) = S_0 - \mu\,t$ with $\mu > 0$.
The phase always decreases (the oscillation
$e^{-i\mu t}$ always winds in the same direction).
Time-reversal ($t \to -t$) would require
$\psi \to \psi^*$, which reverses all momenta
$\mathbf{v} = \nabla S/\meff \to -\mathbf{v}$ and
is not a symmetry of the boundary conditions.

\subsection{The open-system origin}

The irreversibility has a thermodynamic origin.
In the open-system interpretation developed in
Paper~3~\cite{kulangiev_2026c}, the vacuum
condensate is not a closed system: it exchanges
energy with the external reservoir (the Bulk).
The phase advance $dS/dt = -\mu < 0$ corresponds
to the condensate relaxing toward equilibrium with
the Bulk. The forward direction of time is the
direction in which the condensate's free energy
decreases.

This connects three arrows of time:
\begin{itemize}
\item \emph{Thermodynamic arrow}: entropy increases
  toward equilibrium with the Bulk.
\item \emph{Cosmological arrow}: the condensate
  evolves (expands) as it relaxes.
\item \emph{Quantum-mechanical arrow}: the phase
  accumulates monotonically.
\end{itemize}
All three are manifestations of the same underlying
process: the condensate is not in its ground state
with respect to the Bulk, and its evolution toward
equilibrium defines the forward direction.


% ====================================================================
\section{The Hubble Flow as Phase Winding}
\label{sec:hubble}

The Hubble parameter $H_0$ measures the rate at
which the universe expands. In the superfluid vacuum
framework, the expansion is the condensate's
macroscopic flow.

In the Madelung formalism, the macroscopic flow
velocity is the spatial gradient of the phase:
$\mathbf{v} = \nabla S / \meff$. A uniform Hubble
expansion $\mathbf{v} = H_0\,\mathbf{r}$ therefore
requires a spatial phase profile of the form:
\begin{equation}
    S(\mathbf{x})
    = \frac{1}{2}\,\meff\,H_0\,|\mathbf{x}|^2\,.
\label{eq:S_hubble}
\end{equation}
The Hubble parameter is physically realized as the
curvature of the spatial phase field: $H_0 =
\nabla^2 S / (3\,\meff)$. From the Friedmann
equation derived in
Paper~3~\cite{kulangiev_2026c} with $w = -1$:
\begin{equation}
    H_0^2 = \frac{8\pi G}{3}\,\rho_{\mathrm{eff}}\,.
\end{equation}
The expansion of the universe is not the stretching
of a pre-existing spatial metric; it is the active,
continuous winding of the condensate's spatial phase
as the system evolves toward equilibrium with the
Bulk reservoir. In the phase counter picture, $H_0$
is the rate at which new phase cycles are laid down
per unit comoving distance.


% ====================================================================
\section{Temporal Limits}
\label{sec:limits}
% ====================================================================

\subsection{The minimum time interval}

The fastest possible phase oscillation of the
condensate occurs at the maximum chemical potential,
which is bounded by the healing-length
regularization~\cite{kulangiev_2026e}:
\begin{equation}
    \mu_{\max} = \frac{\meff\,c^2}{\heff}\,,
\end{equation}
giving a minimum resolvable time interval:
\begin{equation}
    \boxed{\;\Delta t_{\min}
    = \frac{2\pi}{\mu_{\max}}
    = \frac{2\pi\,\heff}{\meff\,c^2}
    \sim \tP\;}
\label{eq:dt_min}
\end{equation}
Below this, the concept of ``before and after''
loses operational meaning: the condensate cannot
support coherent phase evolution on shorter
timescales. This is not a discretization of time
but a resolution limit of the emergent clock.

\subsection{Time at acoustic horizons}

At the acoustic horizon of a black hole analogue,
the inflow velocity reaches the sound speed:
$v(r_H) = c_s$. The proper time factor becomes:
\begin{equation}
    \frac{d\tau}{dt}
    = \sqrt{1 - v^2/c_s^2}
    \;\xrightarrow{r \to r_H}\; 0\,.
\end{equation}
The phase counter stops. For external observers,
the condensate at the horizon is phase-locked: no
further oscillation cycles accumulate. This
reproduces the infinite gravitational redshift at
the Schwarzschild radius without requiring a
metric singularity---the condensate remains
well-defined and continuous; only the emergent
clock ceases to function.

Inside the horizon ($v > c_s$), the phase
accumulation resumes in a different mode: the
``time'' and ``radial'' coordinates exchange roles
in the acoustic metric, just as in the
Schwarzschild interior. The condensate continues
to evolve, but the internal clock is no longer
synchronized with external observers.

\subsection{No time before the condensate}

In the open-system
interpretation~\cite{kulangiev_2026c}, the vacuum
condensate nucleated from the Bulk reservoir as a
phase transition. Before this nucleation, there was
no condensate, no phase oscillation, and therefore
no emergent time. The question ``what happened
before the Big Bang?'' is ill-posed: time is
generated by the condensate, and asking about
``before'' the condensate is like asking what is
north of the North Pole.

This resolves the initial singularity problem
differently from other approaches. In loop quantum
cosmology~\cite{bojowald_2001}, the singularity is
replaced by a bounce---time continues through the
transition. In the superfluid vacuum, time itself
begins at the nucleation event. There is no
``before'' because the clock did not exist.


% ====================================================================
\section{Discussion}
\label{sec:discussion}
% ====================================================================

\subsection{Relation to other approaches}

The emergent time picture shares features with
several existing proposals while differing in
mechanism:

\emph{Barbour's timeless mechanics}~\cite{barbour_1999}
treats time as emergent from the configuration space
of the universe. Our approach is compatible: the
condensate's configuration space is
$\{\rho(\mathbf{x}), S(\mathbf{x})\}$, and $t$ is
a path parameter through this space. The difference
is that we identify the specific physical mechanism
(phase oscillation) that generates the parameter.

\emph{The thermal time
hypothesis}~\cite{connes_1994} defines time flow
from the modular automorphism of a thermal state.
In our framework, the condensate's thermal
equilibrium with the Bulk provides the thermal
state, and the phase advance is the modular flow.
The approaches may be equivalent at a formal level.

\emph{Volovik's emergent
time}~\cite{volovik_2003} in superfluid helium is
the closest analogue. Our contribution is to
connect this explicitly to the PPN-constrained
acoustic metric, gravitational time dilation,
and the arrow of time through the open-system
boundary condition.

\emph{The Wheeler-DeWitt
equation}~\cite{dewitt_1967} describes a
``frozen'' universe with no time evolution. In the
condensate picture, the resolution is that the
universe IS frozen from the Bulk perspective:
the condensate's 3D configuration exists
timelessly. Time is an internal label generated
by the phase dynamics, not an external parameter
imposed on the system.

\subsection{What is and is not claimed}

We do not claim that time is an illusion or that
change does not occur. The condensate genuinely
evolves: $\rho(\mathbf{x})$ and $S(\mathbf{x})$
change. What we claim is that the ``time
dimension'' perceived by internal observers is
constructed from this evolution, not prior to it.
The 3+1 Lorentzian signature of the acoustic
metric is a derived structure, not a fundamental
one.

We also do not claim to resolve the measurement
problem in quantum mechanics. The phase of the
condensate is a classical field; the transition
to quantum mechanics for its excitations involves
additional structure (the quantum potential $Q$)
that is beyond the scope of this paper.

\subsection{Testable consequences}

The emergent time picture makes three predictions
distinct from fundamental 3+1 spacetime:

(i) \emph{Minimum time resolution} $\Delta t_{\min}
\sim \tP$: this is the same prediction as
Paper~5~\cite{kulangiev_2026e}, derived here from
the phase counter picture rather than the Bohm
potential. The two derivations agree, providing a
consistency check.

(ii) \emph{Modified dispersion at the Planck scale}:
if time is phase accumulation at rate $\mu$, then
at energies approaching $\meff c^2$ the effective
time evolution becomes nonlinear (the oscillation
frequency depends on amplitude). This predicts
energy-dependent propagation speeds for ultra-high-energy
photons, testable via gamma-ray burst
observations~\cite{amelino_2013}.

(iii) \emph{No pre-Big-Bang signatures}: if time
began with the condensate, there should be no
cosmological signatures of a pre-existing epoch
(no bounce imprint in the CMB). This contrasts
with loop quantum cosmology, which predicts
observable bounce signatures at large angular
scales~\cite{bojowald_2001}.


% ====================================================================
\section{Conclusion}
\label{sec:conclusion}
% ====================================================================

We have shown that time in the logarithmic superfluid
vacuum framework is not a fundamental dimension but
an emergent parameter generated by the phase dynamics
of the condensate order parameter.

The chemical potential $\mu(\mathbf{x}) =
\meff\,c_s^2/\heff$ sets the local clock rate.
Gravitational time dilation is the spatial variation
of $\mu$; the twin paradox is path-dependent phase
accumulation; the arrow of time is the irreversibility
of phase advance in an open system exchanging energy
with the Bulk; and the Hubble flow is the spatial
phase winding rate.

The condensate is a three-dimensional object. Its
state is fully specified by two spatial fields:
$\rho(\mathbf{x})$ and $S(\mathbf{x})$. The ``time
dimension'' perceived by phonons (photons) is an
artifact of the acoustic metric's Lorentzian
signature, constructed entirely from these 3D fields.

At the healing length $\xi \sim \lP$, the phase
counter saturates, providing a minimum time
resolution $\Delta t_{\min} \sim \tP$. At acoustic
horizons, the phase locks and the emergent clock
stops. Before the condensate nucleated, there was
no phase oscillation and therefore no time.

The framework unifies the thermodynamic,
cosmological, and quantum-mechanical arrows of time
as manifestations of a single process: the
condensate's relaxation toward equilibrium with the
Bulk reservoir.


% ====================================================================
\section*{Acknowledgments}
% ====================================================================

Claude (Anthropic, claude-opus-4-6) and Gemini (Google, Gemini 3.1 Pro) were used as editorial and computational tools during manuscript preparation. All scientific content, equations, and interpretations are the author's own.


\begin{thebibliography}{25}

\bibitem{dewitt_1967}
B.~S.~DeWitt,
Quantum theory of gravity. I. The canonical theory,
Phys.\ Rev.\ \textbf{160}, 1113 (1967).

\bibitem{isham_1993}
C.~J.~Isham,
Canonical quantum gravity and the problem of time,
in \emph{Integrable Systems, Quantum Groups, and
Quantum Field Theories}, NATO ASI Series
(Kluwer, 1993), pp.\ 157--287.

\bibitem{kuchar_1992}
K.~V.~Kucha\v{r},
Time and interpretations of quantum gravity,
in \emph{Proc.\ 4th Canadian Conference on General
Relativity and Relativistic Astrophysics}
(World Scientific, 1992), pp.\ 211--314.

\bibitem{barbour_1999}
J.~Barbour,
\emph{The End of Time: The Next Revolution in Physics}
(Oxford, 1999).

\bibitem{connes_1994}
A.~Connes and C.~Rovelli,
Von Neumann algebra automorphisms and
time-thermodynamics relation in generally covariant
quantum theories,
Class.\ Quant.\ Grav.\ \textbf{11}, 2899 (1994).

\bibitem{volovik_2003}
G.~E.~Volovik,
\emph{The Universe in a Helium Droplet} (Oxford, 2003).

\bibitem{kulangiev_2026a}
B.~B.~Kulangiev,
Conditions for emergent gravitational light bending
from a logarithmic superfluid vacuum,
Preprint (2026).
\url{https://kulangiev.github.io#paper1}

\bibitem{kulangiev_2026b}
B.~B.~Kulangiev,
Emergent Newtonian Gravity and the Gravitational Feynman-Onsager Relation in the Logarithmic Superfluid Vacuum,
Preprint (2026).
\url{https://kulangiev.github.io#combined}

\bibitem{kulangiev_2026c}
B.~B.~Kulangiev,
Cosmological implications of the logarithmic superfluid
vacuum: Dark energy, the Hubble tension, and a vacuum
self-consistency equation,
Preprint (2026).
\url{https://kulangiev.github.io#paper3}

\bibitem{kulangiev_2026d}
B.~B.~Kulangiev,
Frozen vacuum perturbations as a dark matter analogue
in the logarithmic superfluid vacuum,
Preprint (2026).
\url{https://kulangiev.github.io#paper4}

\bibitem{kulangiev_2026e}
B.~B.~Kulangiev,
Singularity regularization and the emergent Planck scale in the logarithmic superfluid vacuum,
Preprint (2026).
\url{https://kulangiev.github.io#paper5}

\bibitem{visser_1998}
M.~Visser,
Acoustic black holes: Horizons, ergospheres and Hawking
radiation,
Class.\ Quant.\ Grav.\ \textbf{15}, 1767 (1998).

\bibitem{pound_1960}
R.~V.~Pound and G.~A.~Rebka,
Apparent weight of photons,
Phys.\ Rev.\ Lett.\ \textbf{4}, 337 (1960).

\bibitem{bojowald_2001}
M.~Bojowald,
Absence of a singularity in loop quantum cosmology,
Phys.\ Rev.\ Lett.\ \textbf{86}, 5227 (2001).

\bibitem{amelino_2013}
G.~Amelino-Camelia,
Quantum-spacetime phenomenology,
Living Rev.\ Rel.\ \textbf{16}, 5 (2013).

\end{thebibliography}

\end{document}
